% ===== §15 The Praxis of the Two Axes =====
\section{The Praxis of the Two Axes}
\label{p17:sec:praxis}

\noindent
Everything to this point has been diagnosis, ontology, and frame. This section asks what, given all of it, two people are to do---and it must answer under a constraint the whole paper has imposed: that one tends the conditions of love and does not manage love itself (\xref{p17:sec:conditions}). A praxis so constrained cannot be a technique for producing intimacy; the attempt to produce it directly is the category error that decoheres it. What a praxis \emph{can} be is a disciplined tending of conditions, and the discipline has a definite shape, given by the two axes the paper has distinguished: the generation of value, and the suppression of power's freezing. The shape is not symmetrical, and its asymmetry resolves a tension that has shadowed the paper from the start.

\subsection{The resolution of the wu-wei tension}
\label{p17:sub:praxis-wuwei}

A paper whose first claim is that estrangement is necessary, and whose cosmology makes equilibrium death and non-settling life, appears to recommend a quietism: if the passing is necessary and the winter is renewal, is the counsel not simply to refrain, to let be, to do nothing? Yet the paper has also insisted, repeatedly and on empirical grounds (\xref{p17:sub:practices-empirical}), that counter-force does not arise of itself, that power unopposed accumulates and freezes, that the lines of flight must be actively guarded against second-order capture (\xref{p17:sub:balance-tension}). The two seem to pull opposite ways: let be, and resist. The resolution is that they govern \emph{different axes}.

\begin{claim}[Non-action for value, action for power]
The two axes call for opposite practical postures, and conflating them is the root practical error. Toward the \emph{first axis}---the generation of value---the right posture is largely \emph{non-action} (\emph{wu-wei}): good value circulates of itself when its conditions are kept, the loop closes when it is not forced, the spiral lifts when it is not commanded; here the work is to refrain from the forcing that decoheres. Toward the \emph{second axis}---the suppression of power's freezing---the right posture is sustained \emph{action}: power accumulates spontaneously and freezes if unopposed, and the counter-force that prevents it must be perpetually generated. The whole practical stance of the paper is thus coherent: \emph{let be} where generativity is concerned, and \emph{act} where freezing is concerned. Quietism errs by extending non-action to the second axis, where it becomes acquiescence in domination; voluntarism errs by extending action to the first, where it becomes the forcing that kills the bloom.
\end{claim}

\noindent
This is why the paper is neither a quietism nor a manual of control. It is a discipline of \emph{directing} the two postures rightly: knowing, at each moment, whether what faces one is a matter of the first axis (where one refrains) or the second (where one acts). And that knowing requires, before any practice, a diagnosis.

\subsection{The first act is differential diagnosis}
\label{p17:sub:praxis-diagnosis}

Because reparable misalignment and structural shadow present identically at the surface (\xref{p17:sub:symptom-distinction}), and because the two axes demand opposite postures, the first practical act is never an intervention but a \emph{diagnosis}: a determination of what kind of estrangement is in play, and from which layer it issues. The nineteen-fold etiology is, in practice, a differential-diagnostic instrument---a means of asking, of a given passing, whether it is endogenous (issuing from the dyad's own code, prediction, or value-geometry, \xref{p17:sec:encoding}--\xref{p17:sec:phase}) or exogenous (pressed upon the dyad from the field of power, \xref{p17:sub:ir-gradient}), and whether it is the necessary fluctuation of a living system or the leading edge of a freezing.

The stakes of the diagnosis are high in both directions, and misdiagnosis is destructive both ways. To mistake an exogenous estrangement---a passing induced by the pressure of a powerful family, say---for an endogenous one, and to respond with the non-action proper to the first axis, is to leave the dyad undefended against an incursion, to ``let be'' what is in fact an invasion, and so to acquiesce in the breaching of the core. To mistake an endogenous estrangement---the necessary winter of a living bond---for an exogenous threat, and to respond with the action proper to the second axis, is to mount a defence against the bond's own life, to force throughput during replenishment, to attack as a malfunction what is a fluctuation. The differential diagnosis is what assigns each passing to its axis; it is the indispensable first act, and the etiology's deepest practical value is that it makes the diagnosis possible.

\subsection{Practice toward endogenous estrangement: the labour of letting-be}
\label{p17:sub:praxis-endo}

Where the diagnosis returns an endogenous estrangement---one of the internal causes, a passing that issues from the dyad's own code, prediction, or geometry---the governing posture is the non-action of the first axis. But non-action is not nothing; it is a definite labour, and a subtractive one, with concrete forms keyed to the internal causes.

Against the unattributable error of \xref{p17:sec:prediction}, the practice is the \emph{suspension of judgment}: when persistent deviation cannot be attributed, to decline the high-confidence attribution to the other's character (``they have changed,'' ``they no longer care'') in favour of holding the error open---to treat it as possible mistranslation or noise rather than as proven defection. This is the direct practical correlate of causes four through six, and it is subtractive: the work is to \emph{withhold} the premature attribution that would convert a fluctuation into a verdict.

Against the arrested loop of \xref{p17:sec:phase}, the practice is the \emph{refusal to force closure}: to allow the estranged phase to exist without demanding its immediate resolution, to recognise the winter as winter (\xref{p17:sec:rose}) and refrain from the anxious forcing of throughput that prevents replenishment. Here the practice is to identify ``this is a dormant phase'' and to cease applying the pressure---``we must be well again now''---that would disturb the gathering of force.

Against the negative curvature of \xref{p17:sub:phase-forms}, the practice is the \emph{preservation of positive curvature}: to refrain, during the estranged phase, from the irreversible acts---the unrecallable word, the unretractable threat of exit---that deposit fresh scar tissue in the geometry and make the next loop harder to lift. What is to be restrained in the winter is not feeling but \emph{irreversible action}: the bond can survive almost any passing that does not, in its passing, lay down curvature the next return cannot cross.

\subsection{Practice toward exogenous estrangement: the labour of active defence}
\label{p17:sub:praxis-exo}

Where the diagnosis returns an exogenous estrangement---one pressed upon the dyad from the field of power---the posture inverts. Here non-action is acquiescence, and the second axis governs: the work is the active defence of the core's conditions against the incursion, conducted with the filtered IR strategies of \xref{p17:sec:practices}.

Against the withdrawal-threat of the powerful elder (cause 13), the practice is \emph{de-leveraging}: the reduction of the dyad's dependence on resources the external power can withdraw, so that the threat of withdrawal loses its coercive force---the building of an economic and social base independent of the periphery's grant. Against the forced binary that the inseparability of identity makes catastrophic (cause 16), the practice is \emph{hedging}: the refusal to be driven into balancing-against or bandwagoning-with the beloved's family, the maintenance of a third position that defends the dyad's autonomy without demanding the partner break their own bond. Against the borrowed sword (cause 12), the practice is the \emph{boundary-work of self-legislation}: the substitution, for appeals to external authority to adjudicate the bond, of the dyad's own internally legislated commitment---the function the series assigned to the non-binding vow, here serving as the core's defence against a periphery that would otherwise legislate for it.\footnote{On the non-binding vow as internally legislated commitment---the dyad's self-legislation in the absence of an external enforcing authority---see \textcite{paperxiii}; in the present context it operates as the boundary-work that defends the core against exogenous power.} And against the inseparability of identity itself (cause 16), the most difficult practice is \emph{reverse domestication}: not the zero-sum opposition of the powerful elder, which would wound the partner and likely fail, but the patient transformation, over time, of an external power-relation into a relation drawn toward the generative---the slow conversion of the periphery's power-over into something nearer power-to.

\noindent
These two labours---the letting-be toward endogenous estrangement and the active defence toward exogenous---exhaust the response to estrangement as such. But the suppression of freezing is wider than the response to any particular passing; it is a standing task, internal as well as external, and it has a repertoire of its own. To that repertoire, and to the dialectical principle that governs how its internal and external parts are weighted, the section now turns.

\subsection{The internal repertoire of counter-power, and its dialectical weighting}
\label{p17:sub:praxis-inner}

The suppression of freezing requires a repertoire of counter-power practices internal to the dyad---and, first, a principled answer to a question the paper's emphasis has so far left implicit. The practical chapters have weighted the external far more heavily than the internal: most of the active defence just described concerns the periphery. Is this a neglect of internal domination? It is not, and the justification is dialectical (\xref{p17:sec:histmat}).

\begin{claim}[The dialectical weighting of the contradiction]
That the praxis weights the external more than the internal is not an oversight but a judgment about the \emph{principal aspect of the contradiction}. The freezing of power has its principal aspect, in most bonds and most of the time, at the periphery: internal power is dilute and continually metabolised by the generative logic, while the dense, structural, least-metabolised power comes from outside (the external causes are seven of the etiology's sixteen, and the most destructive among them). But the principal and secondary aspects \emph{transform}: when a dyad's interior itself begins to freeze---one party losing subjecthood, roles ossifying, dissent disappearing---the principal aspect shifts inward, and the internal repertoire becomes primary. The differential diagnosis (\xref{p17:sub:praxis-diagnosis}) therefore judges not only endogenous-versus-exogenous cause but the present location of the principal aspect; and the weighting of practice follows the principal contradiction, whose location is itself historical and mobile.
\end{claim}

\noindent
The internal repertoire, deployed when the principal aspect lies within, comprises a recognisable family of practices, each opposing a specific mode of internal freezing, and each connecting to a concept the paper has already established. \emph{The preservation of subjecthood} opposes \emph{dependence}: the retention of independent friendships, independent work, independent paths of growth and judgment, not for the sake of separateness but so that the relation retains more than one source of power---the micro-practice of keeping difference alive against its dissolution (\xref{p17:sub:balance-connection}). \emph{The practice of dissent} opposes \emph{silence}: the maintenance of the capacity to say no, to hold a different view of the future, to reopen a settled rule---dissent functioning as the relation's immune system, since a system in which no contrary voice remains is one in which power accumulates unchecked; this is the internal form of the right to justification (\xref{p17:sub:practices-newer}). \emph{Role fluidity} opposes \emph{ossification}: the refusal to let the division of labour---who earns, who cares, who decides---harden into fixed roles, since fixed roles generate fixed power; not mechanical exchange but maintained mobility. \emph{Meta-communication} opposes \emph{hidden power}: the capacity to discuss not the decision but the \emph{rule} by which decisions are made---to ask who decides, why thus, whether otherwise---dragging into the answerable the power that hides in defaults; this is the practice nearest the paper's theory, a standing audit of the power-structure itself. \emph{Third-space practice} opposes \emph{dyadic monopoly}: the joint service of a generative process beyond the two---a shared work, a project, a thing made together---so that the relation does not revolve solely on the axis of control-and-being-controlled, which is the practical form of co-authorship around a shared vacancy rather than around each other (\xref{p17:sec:asymmetry}). \emph{External connectivity} opposes \emph{the closed system}: friends, community, mentors, peers, functioning not only as resources of life but as resources of \emph{balance}, the civil society of the bond that prevents it from becoming a sealed kingdom---which is why the boundary must be semi-permeable (cause 17), the same permeability that admits the noise admitting also the air. And \emph{re-contracting} opposes \emph{institutional rigidity}: the periodic reopening of the relation's own rules---are we still of this mind? is this rule still right?---so that the rules themselves retain generativity rather than freezing, the limiting contract of the series here kept living by its own renegotiation.

\begin{claim}[The internal counter-power principle]
The internal practices share one principle: counter-power within the dyad does not eliminate power---impossible among symbolic subjects---but continually creates the conditions under which power can be \emph{seen, questioned, and redistributed}. This is the temporal-internal form of the political economy of conditions (\xref{p17:sec:conditions}): like all the paper's practices, it creates conditions rather than producing the end directly. To preserve subjecthood, to keep dissent alive, to fluidify roles, to make rules discussable, to build a third space and an external web, to re-contract: each keeps a channel of power open to view and to revision, and none of them manufactures the love whose conditions they keep.
\end{claim}

\subsection{The dual-register practice, and the category-error caveat}
\label{p17:sub:praxis-dual}

Drawing the two summits into practice, the suppression of freezing and the cultivation of generative asymmetry operate in both registers the paper has held apart. In the \emph{real}, counter-power is the circling of the shared vacancy rather than of each other (\xref{p17:sec:asymmetry}), and the standing symmetry of the right to withdraw---the ontological fact that neither party can be wholly possessed, cause six read now as a safeguard rather than a wound: a subject one cannot finally predict is a subject one cannot finally dominate. In the \emph{symbolic}, counter-power is the repertoire just given---meta-communication, dissent, the right to justification, hedging, the self-legislated vow as the bond's anti-freezing charter, refusing any party the position of sole legislator over the other.

But the whole repertoire stands under one caveat, and the section fixes it before closing, for without it the practices invert into the error the paper most warns against. \emph{Every practice named here belongs to the level of means.} The seven internal practices, the filtered external strategies, the dual-register counter-power---all of them tend the fence; none of them is the bloom, and none can produce it. A couple that executed the entire repertoire mechanically, perfectly, while no real coupling lived beneath it, would be performing the maintenance of conditions over a void---which is exactly forged trust (\xref{p17:sub:symptom-three}, \cite{paperxiii}): the apparatus immaculate, the substance gone. The practices are effective only as the tending of conditions for a generativity that is actually present; turned into a substitute for that generativity, into a technique for manufacturing it, they decohere the very thing they were meant to protect. The praxis is a gardening, not a fabrication; and a gardener who mistook the tending for the growing would tend an empty bed with great precision.

\subsection{Mixed estrangement, and the ledger of time}
\label{p17:sub:praxis-time}

A last realism qualifies the clean division of the foregoing. Real estrangements are seldom purely endogenous or purely exogenous; they are typically entangled (causes 17 and 18), an external pressure inducing an internal withdrawal, a peripheral power exploiting an internal mismatch, so that the differential diagnosis is rarely a single verdict and more often a weighting of intertwined causes. And the bond keeps a ledger (\xref{p17:sub:phase-forms}): each estrangement and each repair deposits curvature, so that the practice cannot be a series of independent interventions but must be understood as acting on a system with memory, in which today's response bends the geometry through which tomorrow's loop will run. There is, accordingly, no once-and-for-all solution, no settlement that would discharge the task---for a settlement that discharged it would be the equilibrium that is death. The possibility of the enriched return---the spiral, the \emph{復} of a renewed coupling---is not secured by any final intervention but kept open, perpetually, at the confluence of two things the praxis can supply: a diagnosis that correctly assigns each passing to its axis, and a response that meets it with the posture that axis demands. Where the diagnosis is right and the posture is right, the conditions of return are kept; and keeping them, perpetually and without the rest that would be death, is the whole of what a praxis of the two axes can do.
